\documentclass[12pt]{article}
\usepackage{graphicx} % Required for inserting images
\usepackage[margin=0.5in]{geometry}
\usepackage{amsmath, amssymb}
\usepackage{mathrsfs}


\usepackage{soul}


% Dr. David Brown Commands
\newcommand{\ds}{\displaystyle}
\newcommand{\R}{\mathbb{R}}
\newcommand{\N}{\mathbb{N}}
\newcommand{\Q}{\mathbb{Q}}
\newcommand{\C}{\mathbb{C}}


% Commands to get script r
\usepackage{calligra}
\DeclareMathAlphabet{\mathcalligra}{T1}{calligra}{m}{n}
\DeclareFontShape{T1}{calligra}{m}{n}{<->s*[2.2]callig15}{}
\newcommand{\scripty}[1]{\ensuremath{\mathcalligra{#1}}}

% Unit commands
\newcommand{\degree}{\text{°}}
\newcommand{\unitSpeed}{\frac{\text{m}}{\text{s}}}
\newcommand{\unitAcceleration}{\frac{\text{m}}{\text{s}^2}}

% Constant commands
\newcommand{\avogadro}{6.022\times10^{23}}

\newcommand{\boltzmannConstK}{1.381\times10^{-23}\frac{\text{J}}{\text{K}}}
\newcommand{\boltzmannConstInEV}{8.617\times10^{-5}\frac{\text{eV}}{\text{K}}}
\newcommand{\planckConst}{6.626\times10^{-34}\text{J}\cdot\text{s}}

\newcommand{\atmP}{1.01325\times10^5\,\text{Pa}}

% Known Value Commands
\newcommand{\electronMass}{9.109\times10^{-31}\,\text{kg}}
\newcommand{\protonMass}{1.673\times10^{-27}\,\text{kg}}

% Commands to easily write partial derivatives
\newcommand{\partDeriv}[2]{\frac{\partial #1}{\partial #2}}
\newcommand{\doublePartDeriv}[2]{\frac{\partial^2 #1}{\partial^2 #2}}


\title{Problem 7.28}
\author{Gabriel Decker}


\begin{document}


\maketitle
\begin{flushleft}


In this problem, we will consider a free Fermi gas in two dimensions.
This gas is confined to a square with $A=L^2$ instead of a cube.
We assume that the particles in the gas are nonrelativistic.\\~\\


\textbf{Part a}\\
If we first treat the system as $T=0$, then we can calculate the Fermi energy for this system.
I've already done something similar in problem 7.22, but this derivation is even more similar to the book's three-dimensional derivation.
Since there are only two directions, the allowed energies are the following.
$$\epsilon=\frac{h^2}{8mL^2}(n_x^2+n_y^2)$$
$$\implies\epsilon=\frac{h^2}{8mA}(n_x^2+n_y^2)$$
If we again consider $n$-space, this time in two-dimensions, we can use polar coordinates to describe each combination of $n_x$ and $n_y$.
We can use the coordinates $n$, or the distance from the origin, and $\theta$, the angle relative to the $n_x$ axis.
This implies that $n_x^2+n_y^2$ is the radius squared of a circle in $n$-space in polar coordinates.
So, we can write the above $\epsilon$ expression as the following.
$$\implies\epsilon=\frac{h^2n^2}{8mA}$$\\~\\

Fermi energy is the max energy state at $T=0$, which occurs at $n=n_\text{max}$ (see my problem 7.22 solution and the book for more explanation).
$$\implies\epsilon_F=\frac{h^2n_\text{max}^2}{8mA}$$\\~\\

We can more helpfully express the energy or fermi energy by converting it to a relation of $N$.
Similar to problem 7.22 and the book's derivation, we can get the number of particles with an amount of energy, or less, by using the area of a circle in $n$-space.
This corresponds with the using the volume of the sphere in the other derivations.
Problem 7.22 and the book derivation specifically does it for $N$ using the Fermi energy's $n_\text{max}$, but this same technique can be used for any energy less than $\epsilon_F$.\\~\\

Since $n_x$ and $n_y$ are positive integers, all the allowed states within the radius $n$ are in the first quadrant.
So, the area we're interested in is one-fourth of the area of the circle with radius $n$ in $n$-space.
So, the total number of particles with energy less than or equal to $\epsilon(n)$, $N_\epsilon$ is the following.
$$N_\epsilon=2\cdot\frac{1}{4}\cdot(\pi n^2)$$
$$\implies N_\epsilon=\frac{\pi n^2}{2}$$
$N$ is simply the $N_{\epsilon_F}$, so
$$\implies N=\frac{\pi n_\text{max}^2}{2}$$
Plugging this into the equation for $\epsilon$, we get the following.
$$\epsilon=\frac{h^2n^2}{8mA}$$
$$\implies\epsilon=\frac{h^2}{4\pi mA}\cdot\frac{\pi n^2}{2}$$
$$\implies\epsilon=\frac{h^2N_\epsilon}{4\pi mA}$$
$$\implies\epsilon_F=\frac{h^2N}{4\pi mA}$$\\~\\

We can also get the total energy $U$, and thus the average energy $U/N$.
$$U=2\int\int\epsilon(n_x, n_y)\,dn_xdn_y$$
We can evaluate this integral using polar coordinates again, with $dn_xdn_y=n\,dnd\theta$.
$$\implies U=2\int_0^{n_\text{max}}\int_0^{2\pi}\epsilon(n)n\,d\theta dn$$
$$\implies U=2\int_0^{n_\text{max}}\epsilon(n)n\,dn\int_0^{\pi/2}\,d\theta$$
$$\implies U=2\int_0^{n_\text{max}}\frac{h^2n^2}{8mA}n\,dn\cdot\frac{\pi}{2}$$
$$\implies U=\pi\frac{h^2}{8mA}\int_0^{n_\text{max}}n^3\,dn$$
$$\implies U=\frac{\pi h^2n_\text{max}^4}{32mA}$$
$$\implies U=\frac{h^2}{8\pi mA}\cdot\frac{\pi^2 n_\text{max}^4}{4}$$
$$\implies U=\frac{h^2N}{8\pi mA}N$$
$$\implies U=\epsilon_FN/2$$\\~\\

Finally, the average energy is the following.
$$\implies U/N=\epsilon_F/2$$\\~\\~\\


\textbf{Part b}\\
We can now get a formula for the density of states.
Like the book, we can start with the integral for $U$.
$$U=2\int_0^{n_\text{max}}\epsilon(n)n\,dn\int_0^{\pi/2}\,d\theta$$
$$\implies U=\pi\int_0^{n_\text{max}}\epsilon n\,dn$$\\~\\

We want to replace the integral by $n$ to an integral by $\epsilon$, so we need to express $n$ in terms of $\epsilon$.
$$\epsilon=\frac{h^2}{8mA}n^2$$
$$\implies n=\sqrt{\frac{8mA}{h^2}}\sqrt{\epsilon}$$
$$\implies dn=\frac{1}{2}\sqrt{\frac{8mA}{h^2}}(\epsilon)^{-1/2}\,d\epsilon$$\\~\\

Plugging this into $U$, we get the following.
$$\implies U=\pi\int_0^{\epsilon_F}\epsilon\cdot\sqrt{\frac{8mA}{h^2}}\sqrt{\epsilon}\cdot\frac{1}{2}\sqrt{\frac{8mA}{h^2}}(\epsilon)^{-1/2}\,d\epsilon$$
$$\implies U=\int_0^{\epsilon_F}\epsilon\cdot\left[\frac{4\pi mA}{h^2}\right]\,d\epsilon$$\\~\\

The quantity in brackets being multiplied by $\epsilon$ is the number of single-particle states per unit energy.
Notice that the bracketed quantity equals $N/\epsilon_F$.
This is the density of states $g(\epsilon)$ function that we want.
Notice that it is independent of $\epsilon$, so $g(\epsilon)\rightarrow g$.
So,
$$g=\frac{4\pi mA}{h^2}$$\\~\\~\\



\textbf{Part d}\\
Since $g(\epsilon)$ is constant in $\epsilon$, we can solve the following integral derived in the book.
$$N=\int_0^\infty g(\epsilon)\frac{1}{e^{(\epsilon-\mu)/kT}+1}\,d\epsilon$$
$$\implies N=g\int_0^\infty\frac{1}{e^{(\epsilon-\mu)/kT}+1}\,d\epsilon$$
Plugging this integral into sympy, we get the following result.
$$\implies N=gkT\ln(e^{\mu/kT} + 1)$$
$$\implies N=\frac{4\pi mA}{h^2}kT\ln(e^{\mu/kT} + 1)$$\\~\\

With this, we can get an equation for $\mu$ in terms of $N$.
$$\implies \ln(e^{\mu/kT}+1)=\frac{h^2}{4\pi mAkT}N$$
$$\implies e^{\mu/kT}+1=\text{exp}\left(\frac{h^2}{4\pi mAkT}N\right)$$
$$\implies e^{\mu/kT}=\text{exp}\left(\frac{h^2}{4\pi mAkT}N\right)-1$$
$$\implies \mu/kT=\ln\left[\text{exp}\left(\frac{h^2}{4\pi mAkT}N\right)-1\right]$$
$$\implies \mu=kT\ln\left[\text{exp}\left(\frac{h^2}{4\pi mAkT}N\right)-1\right]$$
$$\implies \mu=kT\ln\left(e^{N/gkT}-1\right)$$
$$\implies \mu=kT\ln\left(e^{\epsilon_F/kT}-1\right)$$\\~\\

As we vary $T$, we can see the following.
when $kT\gg\epsilon_F$, this implies that the natural log goes negative, so $\mu$ becomes negative.
When $kT\ll\epsilon_F$, the exponential gets very large, so that implies the following.
$$\implies \mu\approx kT\ln\left(e^{\epsilon_F/kT}\right)$$
$$\implies \mu\approx kT\epsilon_F/kT$$
$$\implies \mu\approx \epsilon_F$$
This is the expected behavior.\\~\\



\textbf{Part e}\\
We must now show that when $kT\gg\epsilon_F$, this implies that the chemical potential of this system is the same as that of an ideal gas.
We can start by reminding ourselves that this condition implies that the argument of the exponential is very small, so we can use the taylor approximation $e^x\approx1+x$.
$$\mu=kT\ln\left(e^{\epsilon_F/kT}-1\right)$$
$$\implies \mu=kT\ln\left(1+\frac{\epsilon_F}{kT}-1\right)$$
$$\implies \mu=kT\ln\left(\frac{\epsilon_F}{kT}\right)$$
$$\implies \mu=kT\ln\left(\frac{h^2N}{4\pi mAkT}\right)$$\\~\\

In chapter 6, the statistical approach for ideal gas was investigated.
It gave the following equation for $\mu$.
\begin{equation}
    \mu=-kT\ln\left(\frac{VZ_\text{int}}{Nv_Q}\right)
\end{equation}
where
\begin{equation}
    v_Q=\ell_Q^3=\left(\frac{h}{\sqrt{2\pi mkT}}\right)^3
\end{equation}\\~\\

For our two-dimensional system, we will consider $\ell_Q^3\rightarrow\ell_Q^2$ and $V\rightarrow A$.
Since we have the spin-up and spin-down internal states, $Z_\text{int}=2$.
Let's modify our equation for $\mu$ to make it look like a two-dimensional version of the above.
$$\mu=kT\ln\left(\frac{h^2N}{4\pi mAkT}\right)$$
$$\implies \mu=kT\ln\left(\frac{1}{2}\frac{h^2}{2\pi mkT}\frac{N}{A}\right)$$
$$\implies \mu=kT\ln\left(\frac{N\ell_Q^2}{2A}\right)$$
$$\implies \mu=-kT\ln\left(\frac{2A}{N\ell_Q^2}\right)$$
This matches our expectation.





\end{flushleft}

\end{document}
